\chapter[The Graviton and the Information Ledger]{The Graviton and the Information Ledger \\[0.2cm] \large Unifying Gravity, Light, and Information via Lattice Strain}
\label{app:graviton}

\noindent\textbf{Abstract:}
In the conventional particle description of gravity, the graviton is a spin-two quantum of the gravitational field. Tension Field Theory (TFT) assigns it a different mechanical identity: \textbf{the Graviton is the Longitudinal Strain Wake of a photon propagating across the vacuum lattice.} The photon carries dynamic transverse action through the FCC Soft Floor, while its Tension Column establishes the corresponding longitudinal signature on the SC Hard Floor. Starting from the wavelength-scale loading prescription and the reduced record-energy normalization of TFT, this appendix derives the characteristic graviton energy $E_{\mathrm g}$ and its geometric hierarchy relative to the photon energy $E_\gamma=hf$. The same strain wake supplies the physical interpretation of the information record $N$. Its induced profile tracks the active excitation; its written configuration persists after the excitation has passed, without requiring continued Hard-Floor oscillation. Light, gravitational strain, and persistent information are thereby connected through one dual-scale mechanical architecture.

\section{Introduction: The Mechanical Coupling}
\label{sec:AC_mechanical_coupling}

As established in the Vertical Topology architecture of \hyperref[chap:vertical_topology]{Chapter~\ref*{chap:vertical_topology}}, the Spacetime Gravitational Floor is a stratified mechanical structure. Its dual-phase FCC Active Vacuum is coupled to the phase-locked SC Latent Substrate through the Tension Column.

\begin{itemize}
    \item \textbf{The Photon:} A coherent transverse vibrational wave propagating through the FCC scaffold and its participating interstitial condensate---\textbf{Dynamic Action / The Carrier}.
    \item \textbf{The Graviton:} The corresponding longitudinal strain wake associated with the SC Hard Floor---\textbf{Longitudinal Strain / The Structural Imprint}.
\end{itemize}

Because the Active Vacuum has finite contact stiffness, a photon's coordinated motion involves mechanical stress. In TFT's vertical topology, that active loading establishes a longitudinal substrate signature through the Tension Column.

\noindent\textbf{This signature is the Graviton---the strain wake accompanying the photon.}

The Hard Floor supplies the structural reference rather than another independently propagating finite-energy oscillator. Its limiting incremental rigidity and zero-delay constraint connectivity retain the meanings established in \hyperref[sec:sgf_constitution]{Section~\ref*{sec:sgf_constitution}} and \hyperref[app:physical_wavefunction]{Appendix~\ref*{app:physical_wavefunction}}. \hyperref[def:definition_five]{Definition~V} additionally specifies that a written longitudinal configuration persists after the originating excitation has passed.

The moving signature and the permanent record are therefore related but distinct: the first tracks the present excitation; the second retains its history.

\section[Derivation of the Graviton Energy (Eg)]{Derivation of the Graviton Energy ($E_{\mathrm g}$)}
\label{sec:AC_graviton_energy}

We calculate the characteristic strain-energy scale of the graviton associated with a single photon.

\noindent\textbf{The retained photon regime.}
The wavelength-only expressions below use the uniform, gravity-neutral, long-wavelength photon reduction of \hyperref[app:photon_hamiltonian]{Appendix~\ref*{app:photon_hamiltonian}}, for which $f\simeq c_{\mathrm E}/\lambda$. The loading prescription and record-energy closure specify the reduced normalization; the actual spatial strain profile remains a mode and boundary-value calculation. The approximation signs retain the distinction between this wavelength reduction and the full photon dispersion.

\subsection{The Photon's Energy-Equivalent Mass}

For a monochromatic harmonic excitation with total modal action $\hbar$, \hyperref[app:photon_hamiltonian]{Appendix~\ref*{app:photon_hamiltonian}} gives $E_\gamma=\hbar\omega_\gamma=hf$. Its energy-equivalent loading scale is
\begin{equation}
    m_{\mathrm{eff}}
    =\frac{E_\gamma}{c_{\mathrm E}^{2}}
    =\frac{hf}{c_{\mathrm E}^{2}}.
    \label{eq:AC_effective_mass}
\end{equation}
Here $m_{\mathrm{eff}}$ is the scalar mass equivalent used in the gravitational loading prescription. It is neither an invariant photon rest mass nor the collective mode inertia $M_{\mathrm A}$ defined in \hyperref[app:photon_hamiltonian]{Appendix~\ref*{app:photon_hamiltonian}}.

\noindent\textbf{One photon, one completed action.}
The elementary complete-cycle action remains $h=2\pi\hbar$, as established in \hyperref[sec:X_cycles_and_phase]{Section~\ref*{sec:X_cycles_and_phase}}. Decomposing one photon among participating real modes does not assign a separate elementary quantum to every mode coefficient or lattice contact.

\subsection[The Lattice Strain (alpha G,gamma)]{The Lattice Strain ($\alpha_{\mathrm G,\gamma}$)}

The structural gravitational coupling used in \hyperref[app:substrate_units]{Appendix~\ref*{app:substrate_units}} evaluates a source's loading at a specified distance. For the reduced photon--wake comparison, take that distance to be the photon wavelength:
\begin{equation}
    \alpha_{\mathrm G,\gamma}
    \equiv\frac{Gm_{\mathrm{eff}}}{\lambda c_{\mathrm E}^{2}}.
    \label{eq:AC_photon_loading}
\end{equation}
Substituting Equation~\eqref{eq:AC_effective_mass},
\begin{equation}
    \alpha_{\mathrm G,\gamma}
    =\frac{G\left(hf/c_{\mathrm E}^{2}\right)}
    {\lambda c_{\mathrm E}^{2}}
    =\frac{Ghf}{c_{\mathrm E}^{4}\lambda}.
    \label{eq:AC_frequency_loading}
\end{equation}
Using $f\simeq c_{\mathrm E}/\lambda$ gives
\begin{equation}
    \alpha_{\mathrm G,\gamma}
    \simeq\frac{Gh}{c_{\mathrm E}^{3}\lambda^{2}}.
    \label{eq:AC_wavelength_loading}
\end{equation}

The subscript $\gamma$ identifies the photon-associated loading scale, distinguished from an externally imposed gravitational preload acting on an excitation or detector. The choice $r=\lambda$ fixes the comparison scale; it does not assert that every photon packet is a sphere of radius $\lambda$. The prepared mode, spatial participation, polarization, and interface geometry determine its actual displacement and traction. Outside the stated reduction, Equation~\eqref{eq:AC_frequency_loading} uses the frequency supplied by the actual photon dispersion.

\subsection{The Geometric Substitution: The Hierarchy Law}

The calibrated substrate unit system of \hyperref[app:substrate_units]{Appendix~\ref*{app:substrate_units}} gives
\begin{equation}
    \ell_{\mathrm S}^{2}
    =\frac{\hbar G}{c_{\mathrm E}^{3}}.
    \label{eq:AC_substrate_spacing}
\end{equation}
Using $h=2\pi\hbar$, the photon-associated coupling becomes
\begin{equation}
    \alpha_{\mathrm G,\gamma}
    \simeq2\pi\left(\frac{\ell_{\mathrm S}}{\lambda}\right)^{2}.
    \label{eq:AC_geometric_hierarchy}
\end{equation}

\noindent\textbf{Interpretation:}
Within this wavelength-scale normalization, the gravitational strain associated with a photon is proportional to the square of the ratio between the fundamental SC spacing and the photon wavelength.

\noindent\textbf{The hierarchy is geometric.}
The active wavelength vastly exceeds the structural spacing of the substrate reference. The small coupling follows from that separation of scales, without an additional fitted suppression factor in the reduced normalization.

\subsection[The Graviton Energy (Eg)]{The Graviton Energy ($E_{\mathrm g}$)}

TFT adopts the reduced strain-energy closure that assigns the graviton wake the photon energy multiplied by its longitudinal loading fraction:
\begin{equation}
    E_{\mathrm g}
    \equiv E_\gamma\alpha_{\mathrm G,\gamma}
    \simeq(hf)\left[
    2\pi\left(\frac{\ell_{\mathrm S}}{\lambda}\right)^{2}
    \right].
    \label{eq:AC_record_energy_closure}
\end{equation}
This is the record-energy normalization associated with \hyperref[def:definition_five]{Definition~V} and developed in Section~\ref{sec:AC_information_record}.

Expanding Equation~\eqref{eq:AC_record_energy_closure} in the stated long-wavelength regime gives
\begin{equation}
    E_{\mathrm g}
    \simeq\frac{2\pi h c_{\mathrm E}\ell_{\mathrm S}^{2}}{\lambda^{3}}
    =\frac{h^{2}G}{c_{\mathrm E}^{2}\lambda^{3}}.
    \label{eq:AC_graviton_energy_scale}
\end{equation}

\noindent\textbf{The carrier and its wake scale differently.}
The characteristic graviton energy scales as $\lambda^{-3}$, while the photon energy scales as $\lambda^{-1}$. Their ratio retains the squared geometric suppression in Equation~\eqref{eq:AC_geometric_hierarchy}.

Once the loading and record-energy normalizations are specified, no additional adjustable multiplier enters this reduced hierarchy. Its scale follows from $h$, $G$, $c_{\mathrm E}$, and the selected photon wavelength. The result assigns a characteristic energy to the photon's strain wake, not to an additional independently propagating particle.

\section{Numerical Evaluation and the Gravitational Hierarchy}
\label{sec:AC_numerical_hierarchy}

Consider a standard green visible photon with $\lambda\approx500\,\mathrm{nm}=5.0\times10^{-7}\,\mathrm m$.

\noindent\textbf{The Photon Energy ($E_\gamma$):}
\begin{equation}
    E_\gamma
    \simeq\frac{hc_{\mathrm E}}{\lambda}
    \approx3.9729\times10^{-19}\,\mathrm J
    \approx2.48\,\mathrm{eV}.
    \label{eq:AC_visible_photon_energy}
\end{equation}

\noindent\textbf{The Strain Factor ($\alpha_{\mathrm G,\gamma}$):}
Using $\ell_{\mathrm S}\approx1.6163\times10^{-35}\,\mathrm m$,
\begin{equation}
    \alpha_{\mathrm G,\gamma}
    \simeq2\pi\left(
    \frac{1.6163\times10^{-35}}{5.0\times10^{-7}}
    \right)^{2}
    \approx6.5657\times10^{-57}.
    \label{eq:AC_visible_photon_loading}
\end{equation}

\noindent\textbf{The Graviton Energy ($E_{\mathrm g}$):}
\begin{equation}
    E_{\mathrm g}
    =E_\gamma\alpha_{\mathrm G,\gamma}
    \approx(3.9729\times10^{-19})(6.5657\times10^{-57})\,\mathrm J
    \approx2.6085\times10^{-75}\,\mathrm J.
    \label{eq:AC_visible_graviton_energy}
\end{equation}

\noindent\textbf{The Hierarchy Ratio:}
\begin{equation}
    \frac{E_{\mathrm g}}{E_\gamma}
    =\alpha_{\mathrm G,\gamma}
    \approx6.5657\times10^{-57}
    \sim10^{-56}.
    \label{eq:AC_visible_hierarchy_ratio}
\end{equation}

The separation is enormous, but its origin is explicit: the graviton wake is assigned the longitudinal strain-energy scale associated with a primary transverse excitation. This is the hierarchy of the \textbf{photon--graviton pair} in the stated normalization, rather than a universal numerical ratio between all gravitational and electromagnetic interactions.

\section{The Dynamic Cursor: The Strain Wake}
\label{sec:AC_dynamic_cursor}

\noindent\textbf{The Primary Wave:}
The photon propagates through the Active Vacuum at the $r_{\mathrm D}\sim10^{-26}\,\mathrm m$ structural scale, carrying coordinated transverse motion and contact deformation.

\medskip
\noindent\textbf{The Induced Strain:}
Its Tension Column establishes the associated longitudinal graviton signature at the $\ell_{\mathrm S}\sim10^{-35}\,\mathrm m$ substrate scale. The location of this loading tracks the photon through successive substrate addresses as a dynamic structural cursor.

\medskip
\noindent\textbf{The Retained Trace:}
At a written address, the longitudinal configuration persists after the active excitation has moved on. The cursor identifies the current loading; the retained strain records the passage.

The photon and its induced graviton wake are physically inseparable within TFT's vertical topology. Tracking of the wake does not require substrate pixels to travel with the photon as an ordinary material wave, nor does it require an independent finite-energy oscillator in the bare Hard Floor.

\noindent\textbf{The photon moves. Its graviton wake marks the interaction. The written configuration remains.}

This distinction also separates the permanent substrate record from temporary elastic storage in the Active Vacuum. The interference construction in \hyperref[app:photon_hamiltonian]{Appendix~\ref*{app:photon_hamiltonian}} allows kinetic energy to become momentarily elastic and then return to propagation. Such reversible storage does not, by itself, establish the microscopic transition that writes a permanent Hard-Floor record.

\section[Constitutive Mechanical Definition of the Information Record (N)]{Constitutive Mechanical Definition\\of the Information Record ($N$)}
\label{sec:AC_information_record}

The graviton connects gravitational strain to physical information. The photon is the dynamic carrier of action; its longitudinal wake supplies the structural imprint.

\hyperref[def:definition_five]{Definition~V} expresses information in the TFT strain-bit normalization as the ratio of the carrier's dynamic energy to the static strain energy assigned to a single continuous record:
\begin{equation}
    N=\frac{E_{\mathrm{carrier}}}{E_{\mathrm{record}}}.
    \label{eq:AC_information_ratio}
\end{equation}
For the propagating-wave branch, the adopted closure is $E_{\mathrm{record},\gamma}\equiv E_{\mathrm g}=\alpha_{\mathrm G,\gamma}E_\gamma$. Consequently,
\begin{equation}
    N_\gamma
    =\frac{E_\gamma}{E_\gamma\alpha_{\mathrm G,\gamma}}
    =\frac{1}{\alpha_{\mathrm G,\gamma}}.
    \label{eq:AC_photon_information_ratio}
\end{equation}
The resulting geometric identities retain their branch-specific loading and topological normalizations.

\subsubsection*{A. The Kinematic (Wave) Identity}

Substituting Equation~\eqref{eq:AC_geometric_hierarchy},
\begin{equation}
    N_\gamma
    =\frac{1}{\alpha_{\mathrm G,\gamma}}
    \simeq\frac{1}{2\pi}
    \left(\frac{\lambda}{\ell_{\mathrm S}}\right)^{2}.
    \label{eq:AC_wave_information_identity}
\end{equation}
For the visible-photon example, $N_\gamma\approx1.5231\times10^{56}$. This is the large normalized strain footprint associated with the small longitudinal loading fraction.

\subsubsection*{B. The Inertial (Matter) Identity}

For a coherent standing-wave defect, the information normalization uses its invariant inertial loading $m_{\mathrm i}$ and the substrate mass potential $m_{\mathrm S}=\sqrt{\hbar c_{\mathrm E}/G}$. The matter branch retains the coupling $\alpha_{\mathrm G,\mathrm m}=(m_{\mathrm i}/m_{\mathrm S})^{2}$ and the additional $2\pi$ winding normalization specified in \hyperref[def:definition_five]{Definition~V}.

For the rest-energy comparison $E_{\mathrm{carrier}}=m_{\mathrm i}c_{\mathrm E}^{2}$, consistency with Equation~\eqref{eq:AC_information_ratio} corresponds to the adopted record-energy normalization $E_{\mathrm{record},\mathrm m}=\alpha_{\mathrm G,\mathrm m}E_{\mathrm{carrier}}/(2\pi)$. Hence,
\begin{equation}
    N_{\mathrm m}
    =\frac{2\pi}{\alpha_{\mathrm G,\mathrm m}}
    =2\pi\left(\frac{m_{\mathrm S}}{m_{\mathrm i}}\right)^{2}.
    \label{eq:AC_matter_information_identity}
\end{equation}

\noindent\textbf{The footprint is not the detector count.}
These equations specify the reduced energy normalization of the strain bits in Definition~V. The actual occupied addresses, distinguishable configurations, and repeated-write composition require the microscopic record map. The structural count $N$ is distinct from the active receiving-site count in \hyperref[sec:entanglement_detector_capture]{Section~\ref*{sec:entanglement_detector_capture}}; it is not substituted for the number of detector clicks, a recording threshold, or an outcome probability.

\section{Information Preservation and the Structural Energy Hierarchy}
\label{sec:AC_structural_energy_hierarchy}

The permanence of a written graviton trace belongs to the phase-locked substrate constitution. The accompanying separation between the substrate's structural capacity and the record-energy scale can be evaluated directly.

\noindent\textbf{The Hardware: The Substrate Cell.}
The latent energy-capacity scale assigned to one fundamental SC cell by the substrate unit system is
\begin{equation}
    E_{\mathrm{node}}
    =m_{\mathrm S}c_{\mathrm E}^{2}
    =\sqrt{\frac{\hbar c_{\mathrm E}^{5}}{G}}
    \approx1.956\times10^{9}\,\mathrm J.
    \label{eq:AC_substrate_energy_scale}
\end{equation}
Here $m_{\mathrm S}$ represents latent structural capacity, not the intrinsic inertial mass of an independently oscillating substrate particle.

\medskip
\noindent\textbf{The Software: The Strain Record.}
For the same $500\,\mathrm{nm}$ photon, the reduced record-energy assignment is
\begin{equation}
    E_{\mathrm{record},\gamma}
    \equiv E_{\mathrm g}
    \approx2.6085\times10^{-75}\,\mathrm J.
    \label{eq:AC_record_energy_example}
\end{equation}
Their ratio is
\begin{equation}
    \frac{E_{\mathrm{node}}}{E_{\mathrm{record},\gamma}}
    \approx\frac{1.956\times10^{9}}{2.6085\times10^{-75}}
    \approx7.50\times10^{83}
    \sim10^{84}.
    \label{eq:AC_structural_energy_ratio}
\end{equation}

\noindent\textbf{An energy hierarchy, not a yield calculation.}
Equation~\eqref{eq:AC_structural_energy_ratio} compares two energy scales. A physical yield margin depends on the local stress, geometry, and constitutive response of the writing configuration; it is not determined by this quotient alone.

\medskip
\noindent\textbf{Writing work must be counted once.}
The record-energy assignment is distinct from the work required to write it. \hyperref[app:photon_hamiltonian]{Appendix~\ref*{app:photon_hamiltonian}} already includes the participating kinetic and elastic energy of the photon. Any actual change of stored structural or interaction energy during writing must enter the same complete work and recoil balance, rather than being added again without its corresponding source.

The separation of conservative propagation, receiving work, and registration follows the mechanics of \hyperref[app:matter_vnlw_photon_linear_waveequation]{Appendix~\ref*{app:matter_vnlw_photon_linear_waveequation}}, \hyperref[app:quantum_tunneling]{Appendix~\ref*{app:quantum_tunneling}}, and \hyperref[app:deterministic_measurement_resolution]{Appendix~\ref*{app:deterministic_measurement_resolution}}. The ideal compatibility constraint in \hyperref[app:entanglement]{Appendix~\ref*{app:entanglement}} has no independent net energy uptake. It is not an unmodeled thermal receiver for the writing process.

Thus $E_{\mathrm g}$ is not automatically an energy payment repeated at every microscopic handoff. The identity $N_\gamma E_{\mathrm g}=E_\gamma$ defines the adopted normalization, not an additional energy reservoir.

\medskip
\noindent\textbf{A permanent record and a unitary evolution are different statements.}
Definition~V supplies persistence. A quantum-unitarity correspondence additionally requires the complete evolution and encoding map; it does not follow from the $10^{84}$ energy ratio. The canonical consistency of the retained active dynamics is addressed separately in \hyperref[app:deterministic_measurement_resolution]{Appendix~\ref*{app:deterministic_measurement_resolution}}.

\section{Latent Information Storage: The Holographic Record}
\label{sec:AC_holographic_record}

Within the information premise of TFT, the graviton wake leaves a permanent microscopic strain signature in the deep substrate. Its written longitudinal configuration is the structural record specified by \hyperref[def:definition_five]{Definition~V}.

\noindent\textbf{The excitation passes; the written configuration persists.}
No continuing propagation or oscillation in the Hard Floor is required to maintain that record.

This is the structural meaning of the \textbf{TFT holographic record}: active history is represented in the permanent tension states of the discrete State-0 substrate. The wavelength-squared dependence in Equation~\eqref{eq:AC_wave_information_identity} supplies an area-like normalization. A quantitative encoding capacity additionally requires the distinguishable substrate configurations and their physical writing and probing maps.

\medskip
\noindent\textbf{Persistence is not the same as accessibility.}
The small value of $E_{\mathrm g}$ does not, by itself, establish an absolute threshold of detectability. Readout depends on how an active probe couples to the configuration and produces a measurable response.

As developed in \hyperref[app:entanglement]{Appendix~\ref*{app:entanglement}} and \hyperref[sec:entanglement_decoherence]{Section~\ref*{sec:entanglement_decoherence}}, decoherence may remove access to a shared analyzer response without erasing the underlying substrate history. Observable pair correlations still require the active coupling and detector construction of \hyperref[app:quantum_entanglement]{Appendix~\ref*{app:quantum_entanglement}}.

\subsection{From the Graviton Wake to the Persistent Record}
\label{sec:AC_wake_to_record}

\noindent\textbf{The wake supplies the loading.}
The action-normalized photon field determines the active motion presented to the Tension Column and its participating contacts. The completed action $h$ fixes the elementary normalization; $E_{\mathrm g}$ specifies the adopted strain-energy scale. Neither scalar alone determines the longitudinal pattern written across the substrate addresses.

\medskip
\noindent\textbf{The interface selects the pattern.}
The objective FCC contact law converts relative contact deformation into mechanical traction. Record formation requires the further cross-floor coupling that maps the complete interface history to a retained SC configuration. That configuration must satisfy the substrate's compatibility and constraint conditions; knowing the active traction does not, by itself, supply the writing map.

\medskip
\noindent\textbf{The next write encounters the existing record.}
A rule established for an initially unrecorded substrate does not automatically determine the response of an already written address. Repeated or overlapping interactions require an update law that accounts for the pre-existing longitudinal configuration. Whether successive traces combine additively, modify the local response, or encode their chronological order must follow from that law, not from the single-write normalization.

The same parent mechanics must close the writing, update, and probe interactions, including reciprocal loading, energy, and recoil without counting the photon energy twice. As specified in \hyperref[sec:DMR_structural_limits]{Section~\ref*{sec:DMR_structural_limits}}, the active-coordinate reduction is not analytically continued through its singular endpoint to supply this latent-state dynamics.

\section{The Information--Energy Inverse Law}
\label{sec:AC_information_energy_inverse}

TFT distinguishes structural information from thermodynamic entropy. Entropy characterizes the multiplicity of accessible dynamical configurations under specified macroscopic conditions; structural information concerns the longitudinal configuration retained beneath that evolving active state.

\noindent\textbf{The inverse scaling.}
Within the branch normalizations above, the assigned footprint decreases as the excitation's loading increases. For the long-wavelength photon family and the rest-frame matter branch, respectively,
\[
    N_\gamma\propto\lambda^{2}\propto E_\gamma^{-2},
    \qquad
    N_{\mathrm m}\propto m_{\mathrm i}^{-2}.
\]
Using the substrate energy scale in Equation~\eqref{eq:AC_substrate_energy_scale}, the identities become
\[
    N_\gamma
    \simeq2\pi\left(\frac{E_{\mathrm{node}}}{E_\gamma}\right)^{2},
    \qquad
    N_{\mathrm m}
    =2\pi\left(
    \frac{E_{\mathrm{node}}}{m_{\mathrm i}c_{\mathrm E}^{2}}
    \right)^{2}.
\]
The common inverse-square energy dependence follows algebraically from the adopted wave and matter normalizations, while retaining their different gravitational-coupling conventions.

\medskip
\noindent\textbf{Waves: A Broad Normalized Footprint.}
Within the specified photon family, a lower carrier frequency gives a larger normalized strain footprint. Its actual spatial extent and address pattern still depend on the prepared mode; wavelength alone does not specify the volume of an arbitrary packet.

\medskip
\noindent\textbf{Matter: Concentrated Invariant Loading.}
Increasing the invariant loading of a matter defect reduces its normalized footprint. The value $N_{\mathrm m}=2\pi$ occurs specifically at $m_{\mathrm i}=m_{\mathrm S}$; whether a particular active defect family reaches that scale is a separate phase-boundary question.

These scalings concern the footprint assigned to the current excitation, not a count of every record accumulated along its history. A change in that footprint does not imply erasure of previously written substrate configurations.

\section{Conclusion: Unification via Lattice Mechanics}
\label{sec:AC_conclusion}

TFT connects light, gravitational strain, and information through the vertical structure of the SGF. Under the wavelength-scale loading and record-energy normalization in the retained long-wavelength regime,
\[
    \alpha_{\mathrm G,\gamma}
    \simeq2\pi\left(\frac{\ell_{\mathrm S}}{\lambda}\right)^{2},
    \qquad
    E_{\mathrm g}
    \simeq\frac{h^{2}G}{c_{\mathrm E}^{2}\lambda^{3}},
    \qquad
    N_\gamma=\frac{1}{\alpha_{\mathrm G,\gamma}}.
\]
The coupling, wake energy, and normalized information footprint are linked by the same substrate-to-wavelength geometry. These relations establish the reduced hierarchy; the event-specific writing and readout remain the microscopic completion described in Section~\ref{sec:AC_wake_to_record}.

\noindent\textbf{In TFT, the Graviton is the Strain Wake---not an independent, free-traveling exchange particle. The photon carries the action; the substrate retains the record.}
